\subsection{Uso de \texttt{var}}

\subsubsection{Regla: solo cuando el tipo es evidente}

\textbf{Regla:} Se usa \texttt{var} solo cuando el tipo es evidente en el lado derecho de la asignación. El tipo es evidente únicamente en estos casos:

\begin{itemize}
  \item una construcción con \texttt{new} y el nombre del tipo, con argumentos o con inicializador;
  \item una conversión explícita \texttt{(T)expresion} o el operador \texttt{as T};
  \item una llamada a un método estático que devuelve el mismo tipo que lo declara, como \texttt{Board.Create(...)}, que devuelve \texttt{Board}.
\end{itemize}

En cualquier otro caso se escribe el tipo explícito: tipos integrados (\texttt{int}, \texttt{string}, \texttt{bool} y sus arreglos), aunque se construyan con \texttt{new}, y llamadas a métodos cuyo tipo de retorno no aparece en la expresión. Esta definición coincide con la que aplican las reglas \texttt{IDE0007} e \texttt{IDE0008} con la opción \texttt{csharp\_style\_var\_when\_type\_is\_apparent}, de modo que la compilación la verifica (Microsoft, 2022c).

\textbf{Código correcto:}
\begin{lstlisting}[style=csharpstyle]
var player = new Player("Alice");
var board = Board.Create(BoardSize);
var human = participant as HumanPlayer;
int remainingTurns = CalculateRemainingTurns();
string[] playerNames = new string[MaximumPlayers];
\end{lstlisting}

\textbf{Código incorrecto:}
\begin{lstlisting}[style=csharpstyle]
Player player = new Player("Alice");
Board board = Board.Create(BoardSize);
var remainingTurns = CalculateRemainingTurns();
var playerNames = new string[MaximumPlayers];
\end{lstlisting}
